\section{Exact three-adic depth of the anchor}
\label{sec:three-adic-anchor-depth}

The coefficient size of the physical anchor can be reduced by using its
exact local depth.  The following calculation concerns only the fixed
point and its integral multiples; it is independent of the prime-pattern
and global height hypotheses.

Put
\[
 K=\mathbb Q(\sqrt2),\qquad D=11-6\sqrt2,\qquad
 E:y^2=x^3-Dx,
\]
and
\[
 P=\left(\frac{57}{8}-\frac72\sqrt2,
          -\frac{165}{8}+\frac{417}{32}\sqrt2\right).
\]
The prime three is inert in $K$.  Let $v$ be the valuation on
$K_3=\mathbb Q_3(\sqrt2)$ normalized by $v(3)=1$.
As usual, the valuation of zero is $+\infty$.

\begin{proposition}\label{anchor3:exact-depth}
For every nonzero integer $c$, the point $[c]P$ is finite and has
nonzero abscissa.  Write
\[
             f([c]P)=\frac1{x([c]P)}=a_c+b_c\sqrt2.
\]
Then
\[
 \min\{v_3(a_c),v_3(b_c)\}=
 \begin{cases}
 0,&c\text{ odd},\\
 2+2v_3(c),&c\text{ even}.
 \end{cases}
\]
In the even case the formal parameter $z=-x/y$ satisfies
\[
                  v(z([c]P))=1+v_3(c).
\]
\end{proposition}
\begin{proof}
The displayed equation has good reduction at three: its discriminant is
$64D^3$, which is a unit.  Direct reduction gives
\[
             \overline P=(\sqrt2,0)\in E(\mathbb F_9).
\]
This is a nonidentity point of order two, with unit abscissa.  Hence
every odd multiple of $P$ is finite with unit abscissa.

The ordinate of $P$ has valuation one, since
\[
 y(P)=\frac{3(-220+139\sqrt2)}{32}
\]
and its parenthesized factor has nonzero residue in $\mathbb F_9$.
The numerator of the tangent slope is a unit: modulo three,
$3x(P)^2-D=-D\ne0$.  Thus
\[
 \lambda=\frac{3x(P)^2-D}{2y(P)}
 \quad\hbox{satisfies}\quad v(\lambda)=-1.
\]
The duplication formula now gives
\[
 \begin{split}
 v(x([2]P))&=v(\lambda^2-2x(P))=-2,\\
 v(y([2]P))&=v\bigl(\lambda(x(P)-x([2]P))-y(P)\bigr)=-3.
 \end{split}
\]
There is no cancellation: the summands in the first difference have
valuations $-2,0$, and those in the second have valuations $-3,1$.
Consequently $v(z([2]P))=1$.

For completeness, the needed formal-group valuation is elementary here.
The formal group of the integral Weierstrass equation has integral
coefficients, and negation sends $z$ to $-z$.  Its multiplication series
therefore have the form
\[
                  [m](Z)=mZ+Z^3g_m(Z^2),
                  \qquad g_m(T)\in\mathcal O_{K_3}[[T]].
\]
If $v(z)\ge1$ and $3\nmid m$, the linear term has strictly smaller
valuation than every higher term, so $v([m]z)=v(z)$.
For $m=3$, the linear term has valuation $v(z)+1$, whereas all higher
terms have valuation at least $3v(z)>v(z)+1$.  Thus
\[
                    v([3]z)=v(z)+1.
\]
Write $c=2\cdot3^ru$, where $3\nmid u$.  Repeated application to
$z([2]P)$ yields $v(z([c]P))=1+r$.  In particular, this parameter
is nonzero, so $[c]P$ is not the identity.  On the formal group,
$x=z^{-2}$ times an integral power series with constant term one.
It follows that
\[
                    v(f([c]P))=2+2r.
\]

Finally, $1,\sqrt2$ form an integral basis for the unramified quadratic
extension $K_3/\mathbb Q_3$, with linearly independent residues over
$\mathbb F_3$.  For any rational $a,b$, not both zero, this gives the
exact equality
\[
                  v(a+b\sqrt2)=\min\{v_3(a),v_3(b)\}.
\]
Applying it to the reciprocal abscissas proves both cases.
\end{proof}

\begin{corollary}\label{anchor3:integer-cutoff}
Let $N\ge1$ and $c\ne0$.  Then
\[
 \frac{a_c}{3^N},\frac{b_c}{3^N}\in\mathbb Z_3
 \quad\Longleftrightarrow\quad
 c\text{ is even and }v_3(c)\ge\left\lceil\frac N2\right\rceil-1.
\]
For a prescribed integer $r\ge0$, the smallest positive $N$ forcing
$v_3(c)\ge r$ is $N=2r+1$.  At this value the anchor implies
\[
                             |c|\ge2\cdot3^r.
\]
\end{corollary}
\begin{proof}
The displayed divisibility condition says that the minimum valuation in
Proposition~\ref{anchor3:exact-depth} is at least $N$.  Since $N>0$,
this excludes the odd case.  In the even case it is exactly
$2+2v_3(c)\ge N$.  The stated rounding follows.  When $r\ge1$,
the value $N=2r$ still admits $c=2\cdot3^{r-1}$, so $2r+1$ is
minimal; for $r=0$, it is the smallest positive value by definition.
The final bound follows from the simultaneous divisibility by two and
$3^r$.
\end{proof}

\paragraph{Application to the physical anchor.}
In the merged-anchor theorem, replace its arbitrary constant $C_*$ by
$3^N$ and its coefficient $S_0=C_*^2$ by $3^{2N}$.
The two-adic argument is unchanged because three is a two-adic unit.
Once the root parity and primary divisibility conditions have been
extracted, the three-adic argument recovers exactly the two conditions
in Corollary~\ref{anchor3:integer-cutoff}, with $c=a-1$.
The existing point domain requires $C_S=S-P\ne O$, so $c\ne0$;
this condition must be retained.  If one instead extended $f$ by
$f(O)=0$, the anchor alone would not exclude $c=0$.

Suppose the inherited global argument needs
\[
                     |c|>(80H)^{10},\qquad H<2^h.
\]
It suffices to choose an integer $r$ satisfying
\[
                       2\cdot3^r>(80\cdot2^h)^{10}
\]
and to take
\[
                   N=2r+1,\qquad C_*=3^{2r+1},
                   \qquad S_0=3^{4r+2}.
\]
These substitutions change only fixed integer coefficients.  They add
no coordinate or equation and do not change the polynomial degree.
Completeness is retained by requiring the already freely chosen source
index $c$ to be divisible by $2\cdot3^r$ in the same fixed local
congruence conditions.  This is a fixed refinement of the construction,
not a uniformity assertion with an increasing prime-pattern modulus.

\begin{corollary}\label{anchor3:numerical-constant}
For $h=2093448$, the smallest integer $r$ satisfying
$2\cdot3^r>(80\cdot2^h)^{10}$ is
\[
                         r=13208226.
\]
The corresponding smallest anchor exponent is
\[
 N=26416453,\qquad C_*=3^{26416453},\qquad
 S_0=3^{52832906}.
\]
In particular, a coefficient-height formula $64S_0^2$ specializes to
\[
                              64\cdot3^{105665812}.
\]
The minimality assertion concerns the stated conservative bound
$H<2^{2093448}$ and this exact anchor criterion.
\end{corollary}
\begin{proof}
The threshold for $r$ is equivalently
\[
                     3^r>5^{10}2^{20934519}.
\]
The following four strict integer comparisons are verified by exact
integer exponentiation in the accompanying checker:
\[
 \begin{array}{ll}
 3^{31867}>2^{50508},&
 3^{15288}>5^{10}2^{24207},\\
 3^{15601}<2^{24727},&
 3^{9779}<5^{10}2^{15477}.
 \end{array}
\]
The integer identities
\[
 \begin{split}
 13208226&=414\cdot31867+15288,\\
 20934519&=414\cdot50508+24207
 \end{split}
\]
give the required lower comparison.  Similarly,
\[
 \begin{split}
 13208225&=846\cdot15601+9779,\\
 20934519&=846\cdot24727+15477
 \end{split}
\]
give $3^{13208225}<5^{10}2^{20934519}$ and hence minimality.
The remaining values follow from $N=2r+1$ and $S_0=3^{2N}$.
\end{proof}

The exact checker accompanying this calculation verifies the source
point, its duplication formulas and valuations, and finite examples of
both parity cases using rational arithmetic in $\mathbb Q(\sqrt2)$.
The statement for all indices follows from the formal-group proof,
not from extrapolating those examples.
